% Building 1, Part II. Move systems as rows.
% One file for the whole part. Each chapter begins with a banner and \section*.

\part*{Part II. Move systems as rows}
\phantomsection\addcontentsline{toc}{part}{Part II. Move systems as rows}

%% ==================================================================
%%  Chapter II.1  (the table)
%% ==================================================================
\section*{Chapter II.1. The table of move systems}
\phantomsection\addcontentsline{toc}{section}{\texorpdfstring{Chapter II.1. The table of move systems \quad move systems $\to$ classes}{Chapter II.1. The table of move systems}}
\begin{center}
{\normalsize move systems $\to$ classes \quad$\cdot$\quad from: Part I; Building 0, Chapters 1, 3, 8, 9, 14, II.1 \quad$\cdot$\quad table, blur, class}\\[4pt]
{\small The rows are now move systems. Every row passes the three laws of moves; the table is about what else they pass.}
\end{center}

Building 0, Chapter II.1 again, with move systems for rows. A \emph{move system} is a bag $G$ with $\circ$, a do-nothing $e$ and undos, where $\circ$ is associative: exactly the laws Chapter 1, P1 found for the moves of a bag. By Chapter 1, T3 every finite one is a bag of moves.

\subsection*{Given}

\begin{lines}
\item[\lab{G1}] (Rows.) $\Z$ with $+$ (Building 0, Ch.\ 1); $\Z/6$ (Ch.\ 3); the circle $\mathbb{T}$ (Ch.\ 8); $\Q^\times$, the non-zero rationals with $\times$ (Ch.\ 2); $(\Z/8)^\times = \{1, 3, 5, 7\}$ with $\times$ (Ch.\ 3); $S_3$ (Chapter 1); $D_4$, the moves of $S_4$ that carry a square's corners $1, 2, 3, 4$ (in order round the square) to corners of the square; $GL_2(\R)$, the straight maps of $\R^2$ with $\det \ne 0$ (Building 0, Ch.\ 14); $F_2$ (Chapter 2); $\Z^2$ (Chapter 2, T3).
\item[\lab{G2}] (Facts.) The lines of those chapters, read as yes/no facts.
\end{lines}

\subsection*{Tests}

\begin{lines}
\item[\lab{D1}] $\mathsf{comm}$: $gh = hg$ for all $g, h$.
\item[\lab{D2}] $\mathsf{fin}$: finitely many members.
\item[\lab{D3}] $\mathsf{cyc}$: some $g$ has every member as a power $g^k$, $k$ in $\Z$.
\item[\lab{D4}] $\mathsf{tf}$: no $g \ne e$ has $g^k = e$ for any $k \ge 1$.
\item[\lab{D5}] $\mathsf{div}$: for every $g$ and every $n \ge 1$ there is $h$ with $h^n = g$.
\end{lines}

\subsection*{The table}

\begin{center}
\renewcommand{\arraystretch}{1.25}
\begin{tabular}{l|ccccc}
 & $\mathsf{comm}$ & $\mathsf{fin}$ & $\mathsf{cyc}$ & $\mathsf{tf}$ & $\mathsf{div}$ \\ \hline
$\Z$ & $\checkmark$ & $\cdot$ & $\checkmark$ & $\checkmark$ & $\cdot$ \\
$\Z/6$ & $\checkmark$ & $\checkmark$ & $\checkmark$ & $\cdot$ & $\cdot$ \\
$\mathbb{T}$ & $\checkmark$ & $\cdot$ & $\cdot$ & $\cdot$ & $\checkmark$ \\
$\Q^\times$ & $\checkmark$ & $\cdot$ & $\cdot$ & $\cdot$ & $\cdot$ \\
$(\Z/8)^\times$ & $\checkmark$ & $\checkmark$ & $\cdot$ & $\cdot$ & $\cdot$ \\
$S_3$ & $\cdot$ & $\checkmark$ & $\cdot$ & $\cdot$ & $\cdot$ \\
$D_4$ & $\cdot$ & $\checkmark$ & $\cdot$ & $\cdot$ & $\cdot$ \\
$GL_2(\R)$ & $\cdot$ & $\cdot$ & $\cdot$ & $\cdot$ & $\cdot$ \\
$F_2$ & $\cdot$ & $\cdot$ & $\cdot$ & $\checkmark$ & $\cdot$ \\
$\Z^2$ & $\checkmark$ & $\cdot$ & $\cdot$ & $\checkmark$ & $\cdot$ \\
\end{tabular}
\end{center}
Reasons for the less obvious entries: $\mathbb{T}$ is not $\mathsf{cyc}$ (if $a$ turns a rational amount its powers are finitely many; if not, $a/2$ is no power of $a$) and is $\mathsf{div}$ (Building 0, Chapter 8, T2). $\Q^\times$: $-1$ has $(-1)^2 = 1$; $2$ has no square root. $(\Z/8)^\times$: every member squares to $1$. $GL_2(\R)$: $-I$ squares to $I$; $\operatorname{diag}(-1, 1)$ has no square root ($\det$ of a square is $\ge 0$). $F_2$: $s$ has no square root (exponent sums, Chapter 2, P5); no power of a reduced word $\ne 1$ is $1$.

\subsection*{Blur, and what appears}

\begin{lines}
\item[\lab{B1}] (Building 0, II.1, B1.) For a set $U$ of tests, glue two rows with the same answers on $U$.
\item[\lab{P1}] (Not sharp.) With all five tests, $S_3$ and $D_4$ share a bag. A separating test is the size: $6$ against $8$. Every other pair of rows is told apart.
\item[\lab{P2}] (Classes.) As $T$ runs over the sets of tests, $\operatorname{class}(T)$ takes $11$ values. Old names: $\operatorname{class}(\mathsf{comm})$, \emph{abelian}; $\operatorname{class}(\mathsf{cyc})$, \emph{cyclic} ($\Z$, $\Z/6$); $\operatorname{class}(\mathsf{tf})$, \emph{torsion-free} ($\Z$, $F_2$, $\Z^2$).
\item[\lab{P3}] (Read off the table.) $\mathsf{cyc} \Rightarrow \mathsf{comm}$; \ $\mathsf{div} \Rightarrow \mathsf{comm}$. Never together: $\mathsf{fin}$ and $\mathsf{tf}$; $\mathsf{fin}$ and $\mathsf{div}$; $\mathsf{cyc}$ and $\mathsf{div}$; $\mathsf{tf}$ and $\mathsf{div}$.
\end{lines}

\subsection*{Theorems}

For every move system with more than one member, not only the ten rows.
\begin{lines}
\item[\lab{T1}] $\mathsf{cyc} \Rightarrow \mathsf{comm}$. \ ($g^j g^k = g^{j+k} = g^k g^j$.)
\item[\lab{T2}] $\mathsf{fin}$ and $\mathsf{tf}$ never together. \ (The powers of $g$ repeat, $g^i = g^j$ with $i < j$, so $g^{j-i} = e$.)
\item[\lab{T3}] $\mathsf{fin}$ and $\mathsf{div}$ never together. \ (By T2's reason each $g$ has $g^m = e$ for some $1 \le m \le |G|$; so with $n = |G|!$ every $n$-th power is $e$, and $g \ne e$ has no $n$-th root.)
\item[\lab{T4}] $\mathsf{cyc}$ and $\mathsf{div}$ never together. \ (Finite: T3. Infinite: it matches $\Z$, and $2h = 1$ has no answer there.)
\item[\lab{T5}] (Matching rows.) $\Z/6$ matches $\Z/2 \times \Z/3$, and $(\Z/8)^\times$ matches $\Z/2 \times \Z/2$; $\Z/6$ and $(\Z/8)^\times$ differ already by size, and $\Z/4$ and $(\Z/8)^\times$ by $\mathsf{cyc}$.
\end{lines}

\subsection*{Remarks}

\begin{lines}
\item[\lab{R1}] (Accidents of a small table; Building 0, II.1, R1.) Two lines of P3 die when the right row is added. $\mathsf{div} \Rightarrow \mathsf{comm}$ dies on the rotations of space: every rotation has $n$-th roots (turn by a $1/n$ of the angle about the same axis), and two rotations about different axes need not commute. ``$\mathsf{tf}$ and $\mathsf{div}$ never together'' dies on $\Q$ with $+$. The other four survive every row: T1--T4.
\item[\lab{R2}] (A conjecture from the table, F6.) Every abelian row with $\mathsf{fin}$ in the table matches a product of clocks $\Z/n_1 \times \dots \times \Z/n_k$. This is true for every finite abelian move system; it is proved in the chapter on modules, from a normal form of matrices.
\item[\lab{R3}] (Old names.) Move system: group. $\mathsf{comm}$: abelian. $\mathsf{tf}$: torsion-free. $\mathsf{div}$: divisible. $D_4$: the dihedral group of order $8$. T2: finite groups are torsion groups.
\item[\lab{R4}] (Machine check, \texttt{checks/b1-p2-ch1-check.py}.) Every finite entry by brute force; the $11$ classes (\texttt{b1-p2-ch1-classes.py}); P1; T5 (the two matchings checked).
\end{lines}

\subsection*{Exercises}

\begin{lines}
\item[\lab{E1}] List the members of $D_4$ as moves of the four corners, and show $D_4$ is not $\mathsf{comm}$.
\item[\lab{E2}] Add the row $\Q$ with $+$. How many classes now?
\item[\lab{E3}] Find a test that separates $S_3$ from $D_4$ without counting members.
\item[\lab{E4}] Show $\Z^2$ is not $\mathsf{cyc}$.
\end{lines}

\newpage

%% ==================================================================
%%  Chapter II.2  (cosets)
%% ==================================================================
\section*{Chapter II.2. Blurs that respect one side}
\phantomsection\addcontentsline{toc}{section}{\texorpdfstring{Chapter II.2. Blurs that respect one side \quad a move system $\to$ its cosets}{Chapter II.2. Blurs that respect one side}}
\begin{center}
{\normalsize a move system $\to$ its cosets \quad$\cdot$\quad from: Chapter II.1; Building 0, Chapters 3, II.2 \quad$\cdot$\quad blur, count}\\[4pt]
{\small Ask a blur to respect moving on the left only. Its bags are copies of the bag of $e$, all of one size. Counting them gives Lagrange, and Fermat's little theorem falls out.}
\end{center}

Building 0, II.2 asked which blurs of a ring respect $+$ and $\times$. Here the question is asked of a move system with less: respect $g \circ$ on the left, and nothing on the right. Rule L7 (say what the blur respects) is the whole story of this chapter and the next.

\subsection*{Given}

\begin{lines}
\item[\lab{G1}] (Chapter II.1.) Move systems; the rows.
\item[\lab{G2}] (Building 0, Chapter 3.) $\Z/n$; $(\Z/n)^\times$, the bags with an inverse for $\times$; Chapter 3, T5 there.
\end{lines}

\subsection*{The blur}

On a move system $G$. \ $O$: for each $g$, the move $x \mapsto g x$ (left only).

\begin{lines}
\item[\lab{B1}] (From below.) For a bag $H \subseteq G$: glue $h \sim e$ for each $h$ in $H$.
\item[\lab{B2}] (From above.) For a sub move system $K$ (closed under $\circ$ and undo, holding $e$): the question $x \mapsto$ the bag $xK := \{x k : k \in K\}$.
\end{lines}

\subsection*{Propositions}

\begin{lines}
\item[\lab{P1}] (Survey: every such blur.) A blur of $G$ respecting left moves is fixed by its bag of $e$, which is a sub move system $K$; and $x \sim y \iff x^{-1} y \in K$. Conversely every sub move system gives one. So these blurs and the sub move systems match.
\item[\lab{P2}] (Appears.) B1 for $H$ equals B2 for $K = $ the smallest sub move system holding $H$. Its bags are the $x K$: the \emph{left cosets}.
\item[\lab{P3}] (Survives: moving a bag.) $x \mapsto g x$ carries the bag $yK$ onto the bag $gyK$, one-to-one. \ (Its undo is $x \mapsto g^{-1} x$.) So all bags have the same size as $K$.
\item[\lab{P4}] (Dies: the right side.) In $S_3$ with $K = \{e, (1\,2)\}$: $e \sim (1\,2)$, but $e (1\,3) = (1\,3)$ and $(1\,2)(1\,3) = (1\,3\,2)$ are in different bags, since $(1\,3)^{-1}(1\,3\,2) = (2\,3) \notin K$.
\item[\lab{P5}] (Normal form.) None in general: no rule picks one member of each $xK$. For $\Z$ and $K = n\Z$ the normal forms $0, \dots, n - 1$ of Building 0, Chapter 3 are such a rule.
\end{lines}

\subsection*{Theorems}

\begin{lines}
\item[\lab{T1}] (Counting bags: Lagrange. A count, not a Survives line.) If $G$ is finite, $|G| = |K| \cdot (\text{number of bags})$. \ (P3: the bags are disjoint and each has $|K|$ members.) So $|K|$ divides $|G|$.
\item[\lab{T2}] (The order of a move divides the size.) For $g$ in a finite $G$, the least $m \ge 1$ with $g^m = e$ divides $|G|$. \ ($K = \{e, g, \dots, g^{m-1}\}$ has $m$ members: T1.)
\item[\lab{T3}] ($g^{|G|} = e$.) From T2.
\item[\lab{T4}] (Building 0, Chapter 3, T5, recovered.) For $a$ with $\gcd(a, n) = 1$: $a^{\varphi(n)} \equiv 1 \pmod n$, where $\varphi(n) = |(\Z/n)^\times|$. For $n = p$ prime: $a^{p-1} \equiv 1$. \ (T3 in $(\Z/n)^\times$.)
\item[\lab{T5}] (Dies: ``every divisor is a size''.) $A_4$ has $12$ members and no sub move system with $6$. \ (Machine-checked; classical.) T1 is ``at most these sizes'', never ``these sizes occur''.
\item[\lab{T6}] (Prime size.) If $|G| = p$ is prime, $G$ matches $\Z/p$. \ (Any $g \ne e$ has order dividing $p$, so $p$: its powers are all of $G$.)
\end{lines}

\subsection*{Remarks}

\begin{lines}
\item[\lab{R1}] (What the blur did.) The cosets are the bags of the one blur that respects left moves; their equal size is a Survives line (P3); Lagrange is a count of bags (T1). The method's own split: blur keeps (P3), counting gives existence and sizes.
\item[\lab{R2}] (Respecting both sides.) If the blur also respects right moves, $K$ must be kept by relabelling, $gKg^{-1} = K$, and the bags themselves form a move system. That is Chapter II.3.
\item[\lab{R3}] (Old names.) Left cosets; index $= $ number of bags; T1: Lagrange's theorem; T4: Euler's theorem, Fermat's little theorem; T5: the converse of Lagrange fails ($A_4$).
\item[\lab{R4}] (Machine check, \texttt{checks/b1-p2-ch2-check.py}.) P1 (for every bag $K \ni e$ of $S_3$ and $D_4$: ``$x^{-1}y \in K$'' is a blur exactly when $K$ is a sub move system); P3; P4; T1--T3 on all sub move systems of $S_4$; T4 for $n \le 200$; T5.
\end{lines}

\subsection*{Exercises}

\begin{lines}
\item[\lab{E1}] List the left cosets of $K = \{e, (1\,2)\}$ in $S_3$, and then the right cosets $Kx$. Are they the same bags?
\item[\lab{E2}] Find $3^{100} \bmod 7$ with T4.
\item[\lab{E3}] Which sizes of sub move systems occur in $S_4$? Compare with the divisors of $24$.
\item[\lab{E4}] Show that a move system of size $4$ is $\mathsf{comm}$. (Each $g \ne e$ has order $2$ or $4$.)
\end{lines}

\newpage

%% ==================================================================
%%  Chapter II.3  (normal bags)
%% ==================================================================
\section*{Chapter II.3. Blurs that respect both sides}
\phantomsection\addcontentsline{toc}{section}{\texorpdfstring{Chapter II.3. Blurs that respect both sides \quad a move system $\to$ its quotients}{Chapter II.3. Blurs that respect both sides}}
\begin{center}
{\normalsize a move system $\to$ its quotients \quad$\cdot$\quad from: Chapters 1, 2, II.2 \quad$\cdot$\quad blur}\\[4pt]
{\small Ask the blur to respect moving on both sides. The bag of $e$ must now be kept by relabelling, and the bags form a move system of their own. Every arrow is such a blur followed by an inside.}
\end{center}

The same gluing as Chapter II.2, with one more operation in $O$. Building 0, Chapter 3 is the case $G = \Z$, where left and right are the same and every sub move system will do.

\subsection*{Given}

\begin{lines}
\item[\lab{G1}] (Chapter II.2.) Left-respecting blurs and cosets; T1 (Lagrange).
\item[\lab{G2}] (Chapter 1.) Relabelling $\sigma^\pi$; the bags of relabelling (types); T1 (survey for $S_n$, $n \le 6$).
\item[\lab{G3}] (Chapter 2.) Presentations; P5 (exponent sums).
\end{lines}

\subsection*{The blur}

On a move system $G$. \ $O$: $x \mapsto gx$ and $x \mapsto xg$ for each $g$.

\begin{lines}
\item[\lab{B1}] (From below.) For a bag $H$: glue $h \sim e$ for each $h$ in $H$.
\item[\lab{B2}] (From above.) For a sub move system $N$ with $gNg^{-1} = N$ for every $g$ (\emph{kept by relabelling}): the question $x \mapsto xN$.
\end{lines}

\subsection*{Propositions}

\begin{lines}
\item[\lab{P1}] (Survey.) A blur respecting both sides is a left-respecting blur (Chapter II.2, P1) whose bag of $e$, $N$, is kept by relabelling: $n \sim e$ gives $g n g^{-1} \sim g g^{-1} = e$. Conversely each such $N$ gives one. Then $xN = Nx$: left and right bags agree.
\item[\lab{P2}] (Appears.) B1 for $H$ equals B2 for the smallest $N$ kept by relabelling that holds $H$ (Chapter 2, T2).
\item[\lab{P3}] (Survives: $\circ$ on bags.) $xN \cdot yN := xyN$ does not depend on the members chosen. The bags form a move system $G/N$, and $x \mapsto xN$ is a picture respecting $\circ$.
\item[\lab{P4}] (Dies: ``kept by relabelling'' does not chain.) In $D_4$ (corners $1, 2, 3, 4$ round a square), let $s$ be the flip $(2\,4)$ and $V = \{e, s, (1\,3), (1\,3)(2\,4)\}$. Then $\{e, s\}$ is kept by relabelling in $V$ ($V$ is $\mathsf{comm}$), and $V$ is kept in $D_4$ (it has half the members). But $\{e, s\}$ is not kept in $D_4$: the quarter turn $r = (1\,2\,3\,4)$ gives $r s r^{-1} = (1\,3)$.
\item[\lab{P5}] (Half is always kept.) If $N$ has half the members of $G$, it is kept by relabelling: the two left bags and the two right bags are both $N$ and the rest.
\end{lines}

\subsection*{Theorems}

\begin{lines}
\item[\lab{T1}] (Every arrow is a blur, then an inside. W4.) For an arrow $f : G \to G'$ respecting $\circ$, the blur ``$f(x) = f(y)$'' respects both sides; its bag of $e$ is $N = \{x : f(x) = e\}$; and $G/N \to G'$, $xN \mapsto f(x)$, is an inside. So $G/N$ matches the picture of $f$. \\
Example: $\operatorname{sgn} : S_n \to \{1, -1\}$ gives $S_n / A_n$ matching $\Z/2$ (Chapter 1, P8).
\item[\lab{T2}] (Blurs of a blur.) The two-sided blurs of $G/N$ match the two-sided blurs of $G$ that are blurrier than $N$'s, by $M/N \leftrightarrow M$; and $(G/N)/(M/N)$ matches $G/M$. Blurring twice is blurring once, with the coarser bag of $e$.
\item[\lab{T3}] (A sub and a blur.) For a sub move system $K$ and $N$ kept by relabelling: $KN$ is a sub move system, $K \cap N$ is kept by relabelling in $K$, and $KN/N$ matches $K/(K \cap N)$. \ (T1 for $K \to G/N$.)
\item[\lab{T4}] (Appears: the blur to $\mathsf{comm}$.) Glue $xy \sim yx$ for all $x, y$, respecting both sides. The bags form a $\mathsf{comm}$ move system $G^{\mathrm{ab}}$, and every arrow from $G$ to a $\mathsf{comm}$ move system passes through it, in exactly one way. From above: the questions are all arrows to $\mathsf{comm}$ systems. \\
$S_n^{\mathrm{ab}}$ matches $\Z/2$, by $\operatorname{sgn}$ ($n \ge 2$). \ $F_2^{\mathrm{ab}}$ matches $\Z^2$, by the exponent sums (Chapter 2, P5). \ $D_4^{\mathrm{ab}}$ matches $\Z/2 \times \Z/2$.
\item[\lab{T5}] (Simple: only two such blurs. A count.) $G$ is \emph{simple} if its only two-sided blurs are ``glue nothing'' and ``glue everything''. $\Z/p$ is simple (Chapter II.2, T1). $A_5$ is simple: its relabelling bags (inside $A_5$) have sizes $1, 15, 20, 12, 12$; a bag of $e$ kept by relabelling is a union of these holding the $1$, and its size divides $60$ (Lagrange); the only such unions have $1$ or $60$ members.
\end{lines}

\subsection*{Remarks}

\begin{lines}
\item[\lab{R1}] (L7 is the whole story.) The same gluing $h \sim e$, with left moves in $O$, gives cosets; with both sides, quotients. Nothing else changed.
\item[\lab{R2}] (Simple pieces.) A ladder of blurs $G \to G/N_1 \to \dots \to \{e\}$ in which each rung's blur is simple lists the \emph{atoms} of $G$ (the atoms ladder of the project, rung 2: the list survives, the gluing does not). $A_5$ is the first atom that is not a clock.
\item[\lab{R3}] (Old names.) Kept by relabelling: normal subgroup. $G/N$: quotient group. T1: the first isomorphism theorem. T2: the correspondence and third isomorphism theorems. T3: the second isomorphism theorem. T4: abelianisation, $G/[G, G]$. T5: simple groups; $A_5$ simple. R2: Jordan--H\"older.
\item[\lab{R4}] (Machine check, \texttt{checks/b1-p2-ch3-check.py}.) P1 on all bags of $S_3$, $D_4$ (and kept-by-relabelling sub systems of $S_4$); P3; P4; P5 on $S_4$; T1 for $\operatorname{sgn}$ and $\det$ on $GL_2(\Z/3)$; T4 sizes for $S_3$, $S_4$, $D_4$; T5 (class sizes of $A_5$, all unions).
\end{lines}

\subsection*{Exercises}

\begin{lines}
\item[\lab{E1}] List the sub move systems of $S_3$ kept by relabelling.
\item[\lab{E2}] Show that every sub move system of a $\mathsf{comm}$ system is kept by relabelling. Which line of Building 0, Chapter 3 is this?
\item[\lab{E3}] Find $S_3^{\mathrm{ab}}$ by gluing $(1\,2)(2\,3) \sim (2\,3)(1\,2)$ and closing up.
\item[\lab{E4}] Why can $A_5$ have no arrow onto $\Z/2$?
\end{lines}

\newpage

%% ==================================================================
%%  Chapter II.4  (modules and the Smith form)
%% ==================================================================
\section*{Chapter II.4. Moves on matrices: the Smith form}
\phantomsection\addcontentsline{toc}{section}{\texorpdfstring{Chapter II.4. Moves on matrices: the Smith form \quad relation matrices $\to$ their normal form}{Chapter II.4. Moves on matrices: the Smith form}}
\begin{center}
{\normalsize relation matrices $\to$ their normal form \quad$\cdot$\quad from: Chapters II.1, II.3; Building 0, Chapters 18, II.2, II.5 \quad$\cdot$\quad blur, build}\\[4pt]
{\small Glue a matrix to every matrix reached by moving its rows and columns. The kept questions are the gcds of the minors; the normal form is a diagonal in which each entry goes into the next. Every finite $\mathsf{comm}$ move system is then a product of clocks.}
\end{center}

The first of five chapters with one template: a move system acts on matrices, the blur forgets the moves, kept questions are found, and a normal form shows they are all. Here the moves are on both sides, $A \mapsto PAQ$.

\subsection*{Given}

\begin{lines}
\item[\lab{G1}] (Building 0, Chapter II.5.) $K[x]$; division with remainder by any $d \ne 0$ when $K$ is a field, and in $\Z$; $\det$ over any ring; P10 (equations are carried); P11 (inverse iff $\det$ is a unit); T2: $\Z$ and $K[x]$ pass the same lines.
\item[\lab{G2}] (Building 0, Chapter 18.) $\gcd$, and ``goes into'' as a one-way blur.
\item[\lab{G3}] (Building 0, Chapter II.2.) Blurs of a ring by ideals; quotients $R/I$.
\item[\lab{G4}] (Chapter II.3.) Quotients of move systems; T1.
\end{lines}

\subsection*{Definitions}

\begin{lines}
\item[\lab{D1}] $R$: $\Z$, or $K[x]$ with $K$ a field. Its units: $\pm 1$, or the non-zero constants.
\item[\lab{D2}] An \emph{$R$-module}: a bag $M$ with $+$ (a $\mathsf{comm}$ move system) and scaling by members of $R$, with the laws of Building 0, Chapter 6, P2. $R^n$ is one; so is every $\mathsf{comm}$ move system, with $R = \Z$ ($k \cdot m := m + \dots + m$).
\item[\lab{D3}] (Build, then blur.) For an $m \times n$ matrix $A$ over $R$: $\operatorname{coker} A := R^n$ blurred by gluing each row of $A$ to $0$, respecting $+$ and scaling. Every finite $\mathsf{comm}$ move system $M$ is a $\operatorname{coker} A$: one generator $x_g$ for each member $g$, one relation $x_g + x_h - x_{g+h}$ for each pair (Chapter 2, T1 and Chapter II.3, T1, read with $+$).
\item[\lab{D4}] $d_k(A)$: the $\gcd$ of all $k \times k$ minors of $A$ (made monic, or $\ge 0$); $d_0 := 1$.
\end{lines}

\subsection*{The blur}

On $m \times n$ matrices over $R$. \ $O$: none beyond the moves themselves.

\begin{lines}
\item[\lab{B1}] (From below.) Gluings: $A \sim PAQ$ for $P$, $Q$ square with $\det$ a unit (G1, P11). Equivalently, by row and column moves: swap two rows; add $r$ times one row to another; scale a row by a unit; and the same for columns.
\item[\lab{B2}] (From above.) Questions: $d_1(A), d_2(A), \dots$.
\end{lines}

\subsection*{Propositions}

\begin{lines}
\item[\lab{P1}] (Survives: the $d_k$.) Each $k \times k$ minor of $PA$ is an $R$-combination of $k \times k$ minors of $A$ (an equation in the entries, carried to every $R$: G1, P10). So $d_k(A)$ goes into $d_k(PA)$, and back again by $P^{-1}$: $d_k$ is kept. The same for $AQ$.
\item[\lab{P2}] (Survives: the module.) $\operatorname{coker}(PAQ)$ matches $\operatorname{coker} A$. \ ($P$ only remixes the relations; $Q$ renames the generators.)
\item[\lab{P3}] (Dies: the entries.) $\begin{pmatrix} 2 & 0 \\ 0 & 3 \end{pmatrix} \sim \begin{pmatrix} 1 & 0 \\ 0 & 6 \end{pmatrix}$ over $\Z$: $d_1 = 1$, $d_2 = 6$ for both.
\item[\lab{P4}] (Normal form; at least one, by Euclid.) Every $A$ is glued to a diagonal \[ \operatorname{diag}(e_1, e_2, \dots, e_r, 0, \dots) \] with $e_1 \mid e_2 \mid \dots \mid e_r$, all $\ne 0$. \ (Move a smallest non-zero entry to the corner; divide with remainder along its row and column; if a remainder is left, it is smaller, and repeat; sizes or degrees cannot fall forever (Building 0, Chapter 1, G8). Then clear the row and column, and make the corner go into every other entry by adding a row.)
\item[\lab{P5}] (At most one.) $e_1 e_2 \cdots e_k = d_k(A)$ up to a unit, so $e_k = d_k / d_{k-1}$: fixed by the kept questions.
\item[\lab{P6}] (Appears.) B1 and B2 agree: two matrices are glued exactly when all their $d_k$ agree. \ (P1 one way; P4 and P5 the other.)
\item[\lab{P7}] (Dies over $\Z[x]$.) The $1 \times 2$ matrix $(2 \;\; x)$ over $\Z[x]$ has $d_1 = 1$, but is glued to no $(e \;\; 0)$: the multiples of $2$ and $x$ together are not the multiples of one member (the constant terms of $2a + xb$ are even). Division with remainder was the line that P4 used, and $\Z[x]$ fails it.
\end{lines}

\subsection*{Theorems}

\begin{lines}
\item[\lab{T1}] (Every module given by finitely many generators and relations.) It matches \[ R^{s} \times R/(e_1) \times \dots \times R/(e_r), \qquad e_1 \mid e_2 \mid \dots \mid e_r, \] for exactly one $s$ and one such list of $e_i$ that are not units (up to units). \ (P2, P4, P5.)
\item[\lab{T2}] (Chapter II.1, R2, proved.) Every finite $\mathsf{comm}$ move system matches a product of clocks $\Z/e_1 \times \dots \times \Z/e_r$ with $e_1 \mid \dots \mid e_r$, in one way only. \ ($R = \Z$, $s = 0$.)
\item[\lab{T3}] (Counting them.) The $\mathsf{comm}$ move systems with $n$ members, up to matching, number $\prod_p \mathrm{p}(a_p)$, where $n = \prod p^{a_p}$ and $\mathrm{p}(a)$ is the number of partitions of $a$. For $n = 8$: $3$ ($\Z/8$, $\Z/2 \times \Z/4$, $\Z/2 \times \Z/2 \times \Z/2$). For $n = 72$: $6$. \ (Split each $e_i$ by primes: Building 0, Chapter 3, T6.)
\item[\lab{T4}] (Two clocks, again.) $\Z/m \times \Z/n$ matches $\Z/mn$ exactly when $\gcd(m, n) = 1$: the Smith form of $\operatorname{diag}(m, n)$ is $\operatorname{diag}(\gcd, \operatorname{lcm})$.
\end{lines}

\subsection*{Remarks}

\begin{lines}
\item[\lab{R1}] (The template.) A move system acting on matrices; the blur ``reachable by a move''; kept questions; a normal form proving they are all. Next in Building 1: equivalence over a field (rank), similarity (Jordan), forms (Sylvester), and the cycle type of Chapter 1 was the first case.
\item[\lab{R2}] (Where existence came from.) Once: the Euclid steps of P4, a construction. Uniqueness came from kept questions (P5).
\item[\lab{R3}] (Old names.) $R$-module; $\operatorname{coker} A$: a presented module. B1: equivalence of matrices. $d_k$: determinantal divisors. P4: the Smith normal form; $e_i$: invariant factors. T1: the structure theorem for finitely generated modules over a principal ideal domain (here Euclidean). T2: the fundamental theorem of finite abelian groups. P7: $\Z[x]$ is not a principal ideal domain.
\item[\lab{R4}] (Machine check, \texttt{checks/b1-p2-ch4-check.py}.) P1 ($d_k$ kept under $200$ random moves); P4--P6 against a Smith form routine on random integer matrices; P3; T3 for $n \le 200$ by listing diagonal forms; T4.
\end{lines}

\subsection*{Exercises}

\begin{lines}
\item[\lab{E1}] Find the Smith form of $\begin{pmatrix} 2 & 4 & 4 \\ -6 & 6 & 12 \\ 10 & -4 & -16 \end{pmatrix}$ over $\Z$, and say which $\mathsf{comm}$ move system it presents.
\item[\lab{E2}] List the $\mathsf{comm}$ move systems with $36$ members.
\item[\lab{E3}] Over $\mathbb{Q}[x]$, find the Smith form of $\begin{pmatrix} x & 1 \\ 0 & x \end{pmatrix}$.
\item[\lab{E4}] Show $\Z/4 \not\cong \Z/2 \times \Z/2$ by P5, and by $\mathsf{cyc}$.
\end{lines}
